\section{} % 3

\subsection{} % 3.1

In patch-based machine learning we do segmentation by training a fixed size kernel
of size $k \times k$, which given a patch of some target image of the same size,
outputs typically a single value, which is the predicted class of the central pixel of the patch.

Given some data set $(X, Y)$ of images and labels, we can divide this into a test and a training set.
The training and test set are then some ratio of the total data set, often 80\% and 20\% respectively.
For patch-based machine learning we would ideally have then that $X$ is a set of
image patches from the training images, and $Y$ are segmented patches -- i.e. labels -- for the central pixel of each patch.
In any case the training and test data is some split which contains raw data and labels.

Segmentation methods which do not rely on learning, have the advantages of first of all being
much less computationally expensive (both memory and time),
and typically they are much more humanly interpretable than, say, black box deep learning neural networks.
Non-learning based methods are typically based on assumptions and insights about the data,
and they tend to perform great on simple segmentation tasks.
But for complex segmentation tasks, the scalability of the neural networks, alongside their
non-linearity, allow these to learn much more complex patterns in data.
So albeit at the greatly increased risk of overfitting, learning-based methods 
are typically much more powerful than non-learning based methods for complex segmentation tasks.


\subsection{} % 3.2

By loading in the model as instructed, and using the method \texttt{model.evaluate}
on the test data, we obtain an accuracy of $0.9577$.
This means that the mean accuracy of the predicted labels at each pixel,
across all the images in the test set, is about said value.

We typically report the evaluation scores on the test set, which the model has not been trained on,
to get a sense of the models ability to generalize to unseen data.
\footnote{
	Note that depending on the data set, a simple train/test split is not enough to
	convey any meaningful indication of generalization, since a data set may have a strong
	bias relative to the set of all possible data points.
}
Given this accuracy, we expect the model to be able to do segmentations on unseen data to some extent,
but not perfectly and potentially with even less accuracy.


\subsection{} % 3.3

First we write a function to extract the patches from an image:

\begin{lstlisting}
def extract_patches(
    image: NDArray[np.uint16], k_size: int = 29
) -> NDArray[np.float32]:
    W, H = image.shape
    assert k_size % 2 == 1, "Kernel size must be odd i.e. have a single center."
    pad = k_size // 2
    padded_img = np.pad(
        image, pad, mode="constant", constant_values=0
    ).astype(np.float32)
    # Shape (N, 29, 29, 1) to match the model's expected input
	# Where N = W * H is the total number of pixels in the original image
    patches = np.zeros((W * H, k_size, k_size, 1), dtype=np.float32)
    
    k = 0
    for i in range(H):
        for j in range(W):
            patch = padded_img[i : i + k_size, j : j + k_size] / 255.
            patches[k, :, :, 0] = patch
            k += 1

    return patches
\end{lstlisting}

This function takes a grayscale image and an (uneven) kernel size $k$,
and returns a $k \times k$ patch for each pixel in the original image.
To achieve this we do zero-padding of the original image by $\lfloor k/2 \rfloor$ pixels on each side.
We then simply extract the $k \times k$ patches by iterating over the original image and slicing the padded image accordingly.
We plot the original image and predicted semgentation for a single image in the test set in figure \ref{fig:og_v_pseg}.

\begin{figure}[H]
	\centering
	\includegraphics[width=0.8\textwidth]{ass8/1128_vs_predicted_seg.png}
	\caption{Original image (left) and predicted segmentation (right) for image 1128 of the \texttt{test\_images}.}
	\label{fig:og_v_pseg}
\end{figure}

In this figure \ref{fig:og_v_pseg} the segmentation mask consist of 58 different classes,
which are represented by different colors using the \texttt{matplotlib.pyplot.get\_cmap("nipy\_spectral")} colormap.

\subsection{} % 3.4

%3.4. (1 point) Load the ground truth segmentation for your test image (in test_images/seg).
% Write a function to evaluate the Dice coefficient against your predicted segmen-
% tation.
% Deliverables: Include the code for your function in the report. Use the func-
% tion to report the mean Dice coefficient computed over all classes.

We load in the true segmentation, which we for visual comparisons sake also plot in figure \ref{fig:tseg_v_pseg}.

\begin{figure}[H]
	\centering
	\includegraphics[width=0.8\textwidth]{ass8/true_vs_pred_mask.png}
	\caption{True segmentation (left) and predicted segmentation (right) for image 1128 of the \texttt{test\_images}.}
	\label{fig:tseg_v_pseg}
\end{figure}

In order to calculate the Dice coefficient, we write the following function:

\begin{lstlisting}
def dice_score(
    y_true: NDArray[np.float32],
    y_pred: NDArray[np.float32],
    num_classes: int = 135,
) -> list[Any]:
    dice_scores = []
    for i in range(num_classes):
        true_class = (y_true == i)
        pred_class = (y_pred == i)
        
        intersection = np.logical_and(true_class, pred_class).sum()
        total_pixels = true_class.sum() + pred_class.sum()
        
        # Only compute dice if the class exists in either the truth or the prediction
        if total_pixels > 0:
            dice = (2.0 * intersection) / total_pixels
            dice_scores.append(dice)
            
    return dice_scores
\end{lstlisting}

Which takes the true and predicted segmentation masks of a single image, and computes the Dice coefficient for each class.
This is achieved by iterating over each possible class, creating boolean masks which then allows us to use the 
\texttt{numpy.logical\_and} function
\footnote{
	One could've used numpy broadcasting and multiplied the two together just as well.
	The logical and is more high level readable in conveying the operation performed.
}
to compute the intersection of the true and predicted class masks.
Summing this intersection is possible, due to the underlying 0-s and 1-s representation of boolean values.
Finally we compute the Dice coefficient using the formula $2 \cdot \text{intersection} / \text{total\_pixels}$,
and we only compute this if there are any pixels in either mask, since we are not interested
in 135 scores each time, but rather $max(|X|, |Y|)$ total dice coefficients.
The mean is obtained by \lstinline{np.mean(dice\_scores(y\_true, y\_pred))}.

When running it on the segmentation mask of 1128 along with the predicted segmentation mask of 1128 
- as seen in figure \ref{fig:tseg_v_pseg} - we obtain a mean Dice coefficient of $0.2085$.

This is a bad score, reflecting the clear discrepancy seen between masks in figure \ref{fig:tseg_v_pseg}.

It is reasonable to conclude that the model is capable of achieving high mean accuracy on pixel-wise classifications,
but that this does not necessarily translate to good segmentation performance.
For instance a very dominant class in either data set is the background class,
which is the most common class in both the true and predicted segmentation masks,
making up a large portion of positives and making true positives easier to achieve for the model.

Another conceivable issue with the model, is its limited field of vision.
It can only every look at patches of size $29\times 29$ losing a lot of spatial information.
For instance, there might be some classes that only occur in the bottom right of the brain,
but this spatial information is not available to the model.
This is a crucial limitation of patch-based segmentation models.
However these are mostly speculations, and would require further analysis of the data and results to make concrete deductions.

\subsection{} % 3.5

By making some up-scaling path in the architecture, we could potentially go from predicting a single pixel label
per model evaluation -- requiring $256^2=65536$ model evaluations for a single $256 \times 256 $ image
-- to predicting all $29^2=841$ per model evaluation.
Augmenting the model will significantly strain the memory requirements. Currently the model has
about $1.8$ million parameters, which would be significantly inflated.
However on the flipside we would only need to do $\frac{256^2}{29^2} \approx 78$ model evaluations per image.
Clearly there is a trade-off in terms of memory and time requirements.

Most interestingly, this architectural change would allow the model to take into account
the whole image rather than just patches, which would grant the model the global spatial information
which it currently has no way of leveraging.
That is we completely disregard patch-based inputs to the model, and instead feed in the whole $256 \times 256$ image
to the model, which then outputs a $256 \times 256$ segmentation mask.
Thus predicting via up-scaling paths - potentially with skip connections like U-net -
the model would be much much larger and a single model evaluation would be substantially more memory and time intensive,
however a single model evaluation would give us a complete segmentation mask for the input image.
